\documentclass[11pt,a4paper]{article}
\usepackage[margin=26mm,headheight=15pt]{geometry}
\usepackage[T1]{fontenc}
\usepackage{lmodern,microtype}
\usepackage{amsmath,amssymb,amsthm,mathtools}
\usepackage{booktabs,longtable,array,tabularx}
\usepackage{xcolor,enumitem,fancyhdr,listings}
\usepackage{xurl}
\usepackage[colorlinks=true,linkcolor=ink,citecolor=ink,urlcolor=teal]{hyperref}
\definecolor{ink}{HTML}{17364A}
\definecolor{teal}{HTML}{087E8B}
\definecolor{pale}{HTML}{EFF5F7}
\definecolor{graytext}{HTML}{52616B}
\hypersetup{pdftitle={U211: a decoder-sharing candidate for Con(PA)},
 pdfauthor={Research prepared for Vladimir Reshetnikov},
 pdfsubject={Explicit 211-state binary Turing machine, mathematical reduction and executable verification}}
\newcommand{\PA}{\mathsf{PA}}
\newcommand{\Con}{\operatorname{Con}}
\newcommand{\Halt}{\operatorname{Halt}}
\newcommand{\UU}{\mathsf U}
\newcommand{\Rot}{M^{\mathrm{rot}}_{212}}
\newcommand{\New}{M_{211}}
\newcommand{\Old}{M_{212}}
\newcommand{\act}[3]{#1\,#2\,#3}
\newcommand{\code}[1]{\texttt{#1}}
\newtheorem{theorem}{Theorem}
\newtheorem{lemma}[theorem]{Lemma}
\newtheorem{proposition}[theorem]{Proposition}
\setlength{\parindent}{0pt}
\setlength{\parskip}{5.5pt plus 1pt}
\setlist{nosep,topsep=5pt,leftmargin=1.6em}
\renewcommand{\arraystretch}{1.13}
\pagestyle{fancy}
\fancyhf{}
\fancyhead[L]{\small\color{graytext}U211 / decoder sharing and Con(PA)}
\fancyhead[R]{\small\color{graytext}Research candidate}
\fancyfoot[C]{\small\thepage}
\renewcommand{\headrulewidth}{0.3pt}
\lstset{basicstyle=\small\ttfamily,breaklines=true,columns=fullflexible,
 keepspaces=true,backgroundcolor=\color{pale},frame=single,
 rulecolor=\color{pale},xleftmargin=4pt,xrightmargin=4pt}
\title{\color{ink}\textbf{U211: a decoder-sharing candidate\\for $\Con(\PA)$}\\[7pt]
 \large Nine closed-source reassociations and a predecessor-tag fold}
\author{Research prepared for Vladimir Reshetnikov}
\date{28 October 2026}

\begin{document}
\maketitle
\thispagestyle{empty}
\begin{abstract}
We give an explicit binary, one-tape Turing machine with 211 working states,
one separate halt, and an all-zero initial tape. It is obtained from the
common supplied U212 machine by nine reassociations of closed source factors
followed by one physical decoder-state identification. The reassociation
exposes two selectors whose incoming transitions leave distinct bits immediately
to the left of the head. Existing decoder states read that bit and implement
both selectors, allowing native state 181 to be deleted. The source retains
all 3,909 primitive occurrences and 197 distinct subtrees.

The package contains the complete table, a deterministic builder, an independent
source audit, a four-case whole-tape proof of the final fold, exact blank startup,
and 5,084 arbitrary-natural primitive contracts covering all 422 final rows.
Together with the supplied mathematical calibration, these establish the intended
PA-formalizable equivalence between blank nonhalting and $\Con(\PA)$.
The complete endpoint is not exported as a Lean, Rocq, or MM0 kernel certificate.
The claim is an explicit, reproducible research candidate with mathematical
and executable evidence; no minimum-state or independently ratified record
claim is made.
\end{abstract}

\vspace{7pt}
\colorbox{pale}{\parbox{\dimexpr\linewidth-2\fboxsep\relax}{
\textbf{Result at a glance.} 211 working states; 422 binary rows; start 0;
dense halt 211; no external initializer. Source: 3,909 primitives, 197
distinct subtrees, depth 42, maximum zero run six. Verification: 5,084
symbolic contracts; all rows covered; exact 48,140-step blank startup.}}

\clearpage
\tableofcontents
\clearpage
\section{Result, model, and proof obligations}
\label{sec:result}

The delivered machine $\New$ uses one tape indexed by $\mathbb Z$, alphabet
$\{0,1\}$, blank symbol 0, and one head. Every working-state/symbol pair has
one transition. Dense working labels are $0,\ldots,210$, the start label is
0, and the distinct halt label is 211. Thus there are exactly 211 working
states and 422 transition rows. Initialization is part of the machine.
Every working label is reachable in the directed transition graph; this
does not assert that all labels occur in its blank-tape execution.

For derivations we retain native labels from the supplied compiler:
native 254 is halt, 84 is the primitive dispatcher, and labels need not
be consecutive. The builder checks a bijective renaming to the dense table.
Native and dense tables describe the same machine. The complete dense table
appears in Appendix~\ref{sec:table} and in \path{machine/u211.tm}.

The new part of the argument has three logically separate components.
First, nine changes of source association preserve the full ordinal
register-program control-flow graph. Second, the final literal machine
implements every source primitive on the prescribed whole-tape frame.
Third, its blank startup reaches exactly the frame required by the inherited
proof-search calibration. The final physical fold also has a short independent
local simulation theorem. Combining the new results with the supplied
finite-refutation calibration gives the intended conclusion
\begin{equation}\label{eq:endpoint}
 \PA\vdash
 \neg\Halt(\New,\mathrm{blank})
 \;\leftrightarrow\;\Con(\UU)
 \;\leftrightarrow\;\Con(\PA).
\end{equation}
Here the justification is a conventional mathematical argument with executable
checks of the finite machine and symbolic runtime obligations. Equation
\eqref{eq:endpoint} is not the output of a newly compiled proof-assistant
kernel. Section~\ref{sec:calibration} spells out the inherited dependencies
and the difference between unary $\UU$ and the actual numerical calculus
searched by the machine. No assertion that $\PA$ is consistent, and no
unconditional assertion of nonhalting, is needed for the equivalence.

\section{Comparison of the supplied 212-state proposals}
\label{sec:comparison}

All comparisons use ProveIt commit
\path{9d50bf3df8c9e14941b549945c4f8400172d81f4}, not a moving branch
head~\cite{snapshot}. The seven complete readable archives contain
byte-identical native U212 tables. Their common source is also identical;
the combined archive stores it in its retained direct-parent directory.
The eighth archive, the one with suffix \code{(5)}, is damaged. Fifty early
local members were recovered with uncompressed sizes and CRC values checked.
The recovered builder, parent and receipts specify the same construction
and expected hashes. Its final table is reconstructed from those inputs,
not directly read from a sound final-table member.

Every proposal applies the same reduction to the same direct-transfer U213
parent, whose native SHA-256 begins \code{bb4f5ac12b29d07d}. The common patch is
\begin{align*}
 180/1 &: \act{1}{R}{172}\ \longmapsto\ \act{1}{R}{211},\\
 172/0 &: \act{0}{L}{112}\ \longmapsto\ \act{0}{L}{67},
 \qquad\text{then delete native 112.}
\end{align*}
Consequently, adding the eight final patches together cannot save another
state: they coincide. Comparisons of different U213 ancestors reproduced
inside these packets should not be mistaken for eight different U212
machines. What can usefully be combined is the proof organization,
independent checks, negative examples, and search information.

\begin{table}[htbp]
\centering\small
\begin{tabularx}{\linewidth}{@{}p{26mm}X@{}}
\toprule
Packet & Useful contribution and assessment\\
\midrule
Stacked, no suffix & Direct whole-tape local proof, endpoint audit, independent
checks, and exploratory bypass/scanner searches. Complete U213 ancestry.\\
Stacked (1) & Clear two-stage proof and time bounds; two local checkers;
explicit six-zero scanner counterexample and the helper-211 stacking conflict.\\
Stacked (2) & Independent path and frame implementations, 6,768 local frames,
calibration audit, and a portable replay.\\
Stacked (3) & Explicit intermediate table, independent two-stage audit,
and a rejected circular-helper alias. Particularly useful final-table guards.\\
Stacked (4) & Most explicit bounded-search ledger: direct aliases, deeper
synchronization, bypasses, contextual targets, and direction changes.\\
Stacked (5) & Recovered builder and receipts corroborate the same result;
the supplied archive is not fully replayable because of its damage.\\
ConU & Explicit source-prefix binding and numerical-theory distinctions;
independent frame checks and the $56+g$ runtime/startup interface.\\
Combined & Compact, flattened ancestry, source/calibration inventory,
preserved hashes and startup/seed checks. Best practical baseline for this
package; organization does not make its theorem stronger.\\
\bottomrule
\end{tabularx}
\caption{Eight proposal packets at the pinned snapshot. ``Stacked'' abbreviates
\protect\path{U212_ConPA_Stacked_2026-10-28}; ``ConU'' and ``Combined'' denote the
two archives dated 2026-10-29. All seven complete packets have the same machine.}
\label{tab:comparison}
\end{table}

The common native U212 SHA-256 is
\begin{center}\small
\path{386ce16259d4c97e4af31b3ba494feaff11a0a1a066e0235dddc154b4cea38c4}.
\end{center}
Its dense table has SHA-256
\begin{center}\small
\path{c979e55e6a6f75b08effdf67067cceb460e30a4c28aef9c79797d0a6d32a074e}.
\end{center}
The machine and calibration chosen for the new work are exactly this
common baseline. The complete combined archive and unary-theory archive
are included intact. The other original packets are identified by exact
URLs, byte lengths, Git blob SHA-1 values and archive SHA-256 values in the
comparison manifest; their much larger duplicate ancestry is not copied
into this deliverable.

All supplied U212 endpoints have the same relevant qualification: they
retain mathematical numerical-calculus, compiler, startup and theory
arguments. The separate fixed K215/ZFC verification~\cite{k215pa} does not
certify a changed table or a different consistency predicate.

\section{Construction: enlarge the closed region, then share a selector}
\label{sec:construction}

The earlier tagged-selector proposal could not simply be stacked with U212:
it repurposed native helper 211, which the accepted U212 patch enters from
$180/1$. The successful construction preserves that helper and changes a
larger closed source region. No axiom is removed and no new proof-search
calculus is introduced.

\subsection{Source reassociation and physical decoder rows}

The source is a full binary sequence tree. Leaves are register primitives
with effects \code{CI}, \code{CD}, \code{CB}, \code{CC}, or \code{CH}; a leaf's
binary address identifies its position in the tree. Branch primitives
refer to enclosing left or right ancestors, so arbitrary reassociation
would not be justified. A factor is \emph{closed} when it has no branch
escaping its root and contains no halt primitive. For three closed factors,
the changes
\[
        ((A,B),C)\quad\longleftrightarrow\quad(A,(B,C))
\]
preserve the ordered primitive effects and their resolved ordinal successors.
The checker establishes closure for the exact factors and then verifies
the entire ordinal control-flow graph, including both cases of every
conditional decrement.

A pure physical decoder state reads a bit, writes that same bit, moves
right, and enters its zero or one child. Write its children as an ordered
pair. If $p=(q,c)$ and $q=(a,b)$, a left rotation makes
$p=(a,q)$ and $q=(b,c)$. If $p=(a,q)$ and $q=(b,c)$, the inverse rotation
makes $p=(q,c)$ and $q=(a,b)$. In both cases the inner decoder $q$ must
have precisely the stated incoming row. This is checked over the whole
literal table, not merely the source tree. The builder also rejects any
incoming left-moving runtime edge to either edited decoder.

The source addresses are changed accordingly. For the first direction,
at each affected prefix $s$, the replacements are
\[
       s00u\mapsto s0u,\qquad
       s01u\mapsto s10u,\qquad
       s1u\mapsto s11u;
\]
the reverse direction uses the inverse prefix code. All occurrences of a
shared physical decoder are transformed. The independent auditor reconstructs
these occurrences from literal prefix executions rather than trusting
the discovery program's map. At every stage it checks all 3,908 internal
prefixes, rejecting a write that changes a bit or a read outside the
known prefix and left sentinel. It independently rechecks the full source
graph after the transformation.

\subsection{The nine rotations}

The useful closed region is rooted at native 174 and occupies primitive
ordinals $[2722,3658)$, with zero-based ordinals. It extends the previously
used 774-primitive emitter by one preceding closed 162-primitive factor.
Its fourteen factors, in unchanged order, are
\[
 159,147,176,98,218,168,98,147,146,218,176,90,211,218.
\]
Their primitive lengths are respectively
\[
 162,52,80,26,56,52,26,52,108,56,80,27,103,56,
\]
which sum to 936. Repeated labels denote repeated occurrences of the same
factor. The nine literal rotations are listed in Table~\ref{tab:rotations}.

\begin{table}[htbp]
\centering
\begin{tabular}{rrrrl}
\toprule
Step & Outer & Inner & Side & Source reassociation\\
\midrule
1&215&204&1&$(A,(B,C))\to((A,B),C)$\\
2&217&181&1&$(A,(B,C))\to((A,B),C)$\\
3&185&217&0&$((A,B),C)\to(A,(B,C))$\\
4&174&185&1&$(A,(B,C))\to((A,B),C)$\\
5&129&158&0&$((A,B),C)\to(A,(B,C))$\\
6&217&215&1&$(A,(B,C))\to((A,B),C)$\\
7&204&214&0&$((A,B),C)\to(A,(B,C))$\\
8&215&129&0&$((A,B),C)\to(A,(B,C))$\\
9&216&220&1&$(A,(B,C))\to((A,B),C)$\\
\bottomrule
\end{tabular}
\caption{The deterministic recipe. Side is the old edge from outer to inner.
All labels are native. Exact occurrence prefixes are in
\protect\path{inputs/rotation_paths.json}.}
\label{tab:rotations}
\end{table}

One way to display the final association, using $(x,y)$ for source
sequencing, is
\begin{align*}
 F&=(211,218), & E&=(218,176),\\
 D&=(146,(E,90)), & C&=(147,D),\\
 B&=((218,(168,98)),C), & A&=(98,B),\\
 \text{root }174&=((159,(147,176)),(A,F)).
\end{align*}
In the physical table, $B$ uses helper 129; its left branch through 158
begins with 218 and its right branch through 204 begins with 147. The
factor represented by 211 remains intact.

After all nine rotations, the source still has 3,909 primitives and 197
distinct subtrees, maximum depth 42 and maximum zero run six. The reduction
does not improve the subtree count: it makes the existing allocation more
shareable. Intermediate rotated tables are construction stages only. Some
intermediate address shapes are unsuitable for the original scanner; the
full runtime theorem is checked for the final table, not asserted for
every intermediate stage.

\section{The predecessor-tag fold}
\label{sec:fold}

Let $\Rot$ denote the 212-state table after the nine rotations. Its relevant
rows are:
\begin{center}
\begin{tabular}{rcc}
\toprule
Native state & Read 0 & Read 1\\
\midrule
220 & $\act{0}{R}{218}$ & $\act{1}{R}{176}$\\
181 & $\act{0}{R}{147}$ & $\act{1}{R}{176}$\\
129 & $\act{0}{R}{158}$ & $\act{1}{R}{204}$\\
158 & $\act{0}{R}{218}$ & unchanged\\
204 & $\act{0}{R}{147}$ & unchanged\\
\bottomrule
\end{tabular}
\end{center}
The complete incoming inventories, including all runtime rows, are
\[
 \operatorname{In}(220)=\{216/0:\act{0}{R}{220}\},\qquad
 \operatorname{In}(181)=\{185/1:\act{1}{R}{181}\}.
\]
Neither state is initial. Thus whenever $\Rot$ enters 220, the cell immediately
left of its head is 0; whenever it enters 181, that cell is 1. The previous
transition wrote the identifying bit, so no condition on earlier tape
contents is needed.

The final fold performs exactly three operations:
\begin{enumerate}
\item Redirect the target of $185/1$ from 181 to 220.
\item Replace $220/0$ by $\act{0}{L}{129}$.
\item Delete the two outgoing rows of 181.
\end{enumerate}
All other retained rows agree with $\Rot$, up to the target identification
$\rho(181)=220$, $\rho(q)=q$ otherwise. In particular, the helpers are checked
in the final edited table, and are distinct from both merged states.

\begin{lemma}[Four local identities]\label{lem:four}
Let the old head position be $h$, with left-neighbor tag 0 at old state 220
or tag 1 at old state 181. One old transition is simulated by one or three
new transitions, with the same whole tape and head position and control
mapped by $\rho$. The exterior tape is arbitrary and is not read.
\end{lemma}
\begin{proof}
The two possible current symbols give the four cases in
Table~\ref{tab:local}. If the current symbol is 1, both old states have
the identical row to 176, which is retained. If it is 0, the new 220 row
moves left to 129. That helper reads and preserves the tag, returning right
to 158 or 204. The current zero then selects the required old target 218
or 147. Every write preserves the symbol read. Only cells $h-1$ and $h$
are read; the head ends at $h+1$, which is not read inside the block.
No displayed intermediate state is halt.
\end{proof}

\begin{table}[htbp]
\centering\small
\begin{tabular}{rrrlc}
\toprule
Old state & Tag & Read & New state/head path, relative to $h$ & Steps\\
\midrule
220&0&0&$220@0\to129@{-1}\to158@0\to218@1$&3\\
181&1&0&$220@0\to129@{-1}\to204@0\to147@1$&3\\
220&0&1&$220@0\to176@1$&1\\
181&1&1&$220@0\to176@1$&1\\
\bottomrule
\end{tabular}
\caption{Exhaustive local cases. The old block has one step in every row.}
\label{tab:local}
\end{table}

\begin{theorem}[Final-fold equivalence]\label{thm:fold}
For any common initial tape, from start state 0 and a common head position,
$\Rot$ halts if and only if $\New$ halts. If a common halting run has
$T_{\mathrm{rot}}$ old steps and $k$ of them read zero at old 220 or 181,
then
\[
 T_{211}=T_{\mathrm{rot}}+2k,\qquad
 T_{\mathrm{rot}}\leq T_{211}\leq3T_{\mathrm{rot}}.
\]
\end{theorem}
\begin{proof}
At time zero neither tagged state is active. The complete incoming lists
establish the required tag at every subsequent old occurrence. For every
old step outside the two tagged states, the corresponding new row commutes
with $\rho$. There are 420 such rows, all checked explicitly. The common
one branch and Lemma~\ref{lem:four} handle the remaining rows.

Induction on old time therefore embeds the old execution into the unique
new execution at block boundaries, with identical tape and head and
projected control. Every block is finite, has length one or three, and
has no internal halt. If the old run halts, the corresponding new boundary
halts. If the old run is infinite, its block times grow without bound and
describe a nonhalting new run; a new halt cannot occur inside a block.
This proves reflection as well as preservation. Summing block lengths
gives the exact time identity and bounds.
\end{proof}

The arbitrary-tape statement is an ordinary mathematical local simulation.
For its PA formulation one codes finite configurations and finite tape
support; blank-tape equivalence and the finite trace induction are arithmetic.
No second-order quantification over infinite tapes is required. The bound
belongs only to the final fold $\Rot\to\New$. The preceding reassociations
change source addresses, so neither this bound nor arbitrary-common-tape
equivalence is asserted for the whole path from the original $\Old$.

The finite structural verifier \path{code/verify_fold.py} checks the
incoming inventories, all ordinary rows, the final helper rows, and the
four paths over partial tapes with an unread opaque exterior. Twelve
intentional corruptions are rejected. The verifier uses explicit checks
and passes also under Python \code{-O}; it does not turn off the obligations
when assertions are disabled.

\section{Whole-tape runtime and blank startup}
\label{sec:runtime}

A source reassociation changes instruction addresses even when it preserves
all ordinal control-flow edges. It therefore requires a fresh connection to
the literal binary machine. We establish that connection by replaying the
inherited symbolic engine~\cite{proposal1} on the final U211 table and every
occurrence in its new source. The engine is preserved unchanged in
\texttt{code/symbolic\_runtime.py}; the new portable driver is
\texttt{code/verify\_runtime.py}. The resulting certificate uses explicit
induction rules for unbounded tape segments. It is executable mathematical
evidence, rather than a proof-assistant kernel certificate.

\paragraph{The quantified frame.}
For a natural number $v$, put
\[
 B(v)=0\,1^{v+1},\qquad C(v)=1\,0\,1^v.
\]
For a finite list $A$, write $B(A)$ and $C(A)$ for the corresponding
concatenations. At a primitive boundary the control is native state~84,
and the head points to the first bit of the instruction address $p$.
Relative to that head, the whole tape has the form
\[
 \cdots 0\,1\;p\;0^{56+g-|p|}\,1\;
 B(v_0)\cdots B(v_8)B(a_0)B(A)\;0\cdots .
\]
Here $g\in\mathbb N$ is an arbitrary extra gap, $a_0\in\mathbb N$ is an
arbitrary first spare value, and $A$ is an arbitrary finite spare tail,
including the empty list. Thus the complete spare list is nonempty.
The nine logical values $x_i$ are arbitrary naturals, arranged physically
by
\[
 \pi=(6,2,0,4,3,7,8,5,1),\qquad v_{\pi(i)}=x_i.
\]
Each guarded decrement is checked separately at zero and at $x_i=u+1$.
The final source has maximum address length~42 and maximum consecutive-zero
run~6; its seven-zero origin detector and minimum header distance~56 remain
compatible with this encoding. Both infinite zero exteriors are represented.
Unmentioned cells are never silently replaced by a finite sample of zeros.

The engine stores run lengths, elapsed time, and displacement as affine
expressions with thirteen slots: a constant, the gap, nine primary values,
the first spare value, and the width of $B(A)$. The last slot denotes the
actual summed width of the represented list; it is not a second, unrelated
choice of tape length. Coefficientwise nonnegativity supplies the elementary
inequalities needed by the following rules.

\paragraph{Four sound symbolic rules.}
The proof construction uses only these operations.
\begin{enumerate}
\item \emph{Literal transition.} A row is read from the final transition
table. A literal read of a run is permitted only when its constant length
term is at least one and all coefficients are nonnegative. Removing its
first symbol, inserting the written symbol, and updating control and head
therefore implement that transition for every represented valuation. A
possibly empty run or a read into an opaque context causes rejection.

\item \emph{Unary scan induction.} If the row preserves the control state
and moves in the current traversal direction, an entire homogeneous run
can be consumed at once. For run length $n$, this abbreviates exactly $n$
applications of the same row. Induction on $n$, including zero, proves the
rule; its time and signed displacement are recorded as $n$ and $\pm n$.
It places no upper bound on a register value.

\item \emph{Finite-list induction.} Before using any list traversal, the
checker rederives its one-block law from the final table, with arbitrary
natural $v$ and opaque contexts on both sides. The four native entry states
are 64, 0, 40, and 97: respectively rightward $B$-to-$B$, rightward
$B$-to-$C$, leftward reversed $B$-to-$B$, and leftward reversed
$C$-to-$B$. Each traversal returns to its entry state after exactly $v+2$
steps and traverses that width. The empty-list case and concatenation of
one more checked block prove the corresponding law for every finite list.

\item \emph{Word normalization.} Adjacent equal-symbol runs are joined,
runs of identically zero length are removed, and the infinite zero exterior absorbs
neighboring finite zero runs. The additional identity
\[
 C(A)1=1B(A)
\]
and its reversal follow from the one-block identity and induction on $A$.
These operations preserve the represented whole word, including its
position relative to the head.
\end{enumerate}

\paragraph{Primitive contracts and positive progress.}
For each source occurrence, the expected endpoint is calculated from its
instruction: increment, the appropriate decrement outcome, an ancestor
cut for a control instruction, or the halt label. Address overflow and the
resulting head displacement are included explicitly. A positive decrement
increases the header gap; an increment appends the prescribed spare block.
For every nonhalting primitive, the resulting control, complete tape, and
head displacement must equal the expected configuration exactly. A halt
case must reach native~254. The duration must be an affine expression with
nonnegative coefficients and a strictly positive constant term.
Consequently every represented source step takes finite positive time.
The bound of 4000 symbolic inferences per attempted contract limits proof
construction, not the natural values or list lengths covered by a successful
contract.

The completed run proves all \textbf{5084 cases}: one for each of 3909
primitive occurrences, together with an additional case for each of the
1175 decrements. It uses \textbf{1,280,108 symbolic inferences}, covers all
\textbf{422 literal transition rows}, and obtains minimum duration constant
\textbf{5}. Transition coverage is a supplementary diagnostic; the quantified
contracts and their induction rules provide the substantive evidence.
The inherited assertion-based engine deliberately rejects optimized Python,
so the portable driver preserves that restriction.

\paragraph{Exact blank initialization.}
A separate sparse-tape replay checks all \textbf{48,140} transitions of the
blank startup. Its entire trace agrees with the U212 baseline~\cite{combined};
no changed retained row is visited. The endpoint is native~84 at head
position $-2$, with source ordinal~136 and address
$\texttt{001100110101010}$. Its header distance is 56, its logical registers
are
\[
 (1,0,13,0,0,1,0,0,1),
\]
and its spare list consists of 43 ones. The complete finite support contains
120 one-cells and agrees with an independently reconstructed frame, which
also fixes the zero exterior. This is an exact finite initialization proof.
Together with the new source's identical ordinal control flow and the
quantified runtime contracts, it transports the baseline's source behavior
to the released literal machine. The inherited safe-suffix and logical
calibration arguments are identified separately in
Section~\ref{sec:calibration}; the checks here do not assert consistency
or unconditional nonhalting.

\section{The inherited consistency calibration}
\label{sec:calibration}

The new machine reduction preserves a specified arithmetic proof search.
Its consistency endpoint combines the new source, runtime, and startup
arguments with the mathematical calibration retained in the supplied U212
packet~\cite{combined}. This section identifies those dependencies. It does
not substitute Con(PA) into the verified theorem about the different K215
ZFC machine~\cite{k215pa}.

\subsection{Unary U and the numerical theory actually searched}

The theory $\UU$ uses classical equality-free first-order logic over a
nonempty domain, with membership as its only nonlogical symbol. For each
formula $\varphi(x)$ with no other free variable, its sole schema is
\[
 \begin{gathered}
 \forall a\,\exists b\,
 \left[
   \forall x\,(x\in b\to\varphi(x))\ \land\
   \forall x\,(x\subseteq a\to(\varphi(x)\to x\in b))
 \right],\\[4pt]
 \text{where }\quad
 x\subseteq a\;:\!\Longleftrightarrow\;
     \forall z\,(z\in x\to z\in a).
 \end{gathered}
\]
Bound variables are freshened. In particular, the source $a$ cannot be
a free matrix parameter. The upper and lower clauses give weak covering;
they do not assert ambient extensionality or an exact power filter.

The supplied unary report~\cite{unary} gives fixed interpretations in both
directions and uniform finite-refutation translations. For the upper
interpretation, PA interprets membership as binary-bit membership.
Inclusion of the one-bits of $x$ among those of $a$ implies $x\leq a$;
formula-specific finite-comprehension induction therefore constructs the
required cover. This is not a uniform truth predicate.

The lower interpretation works with numerical representatives modulo
coextensionality, using saturation wherever membership needs congruence.
Local predecessor facts first give induction without assuming total
successors. Comparison and least numerical failures then yield bounded
numerical comprehension, followed by successors, numerical powers,
pair and graph codes, and addition and multiplication. Parameters are
bound inside legal unary matrices. The resulting proof translations are
primitive recursive, with correctness proved mathematically to be
formalizable in PA. Thus the inherited conclusion is specifically
\[
                  \PA\vdash\Con(\UU)\leftrightarrow\Con(\PA),
\]
not merely an external semantic equiconsistency assertion.

The literal machine searches the retained numerical calculus
$\mathcal D_U$. It uses truth code $0$, falsity code $1$, seven logical
rows, a dominated propositional row, modus ponens, generalization, and
the gated U214 covering bundle. Its machine-facing cover allows additional
matrix parameters and is calibrated through unrestricted covering with
global irreflexivity. The naming conventions, equality interpretation,
small-code aliases, and coincident source/output labels are handled by
the retained U214/U218/U222 proofs~\cite{combined}. These presentations
have finite-refutation translations; their literal axiom sets are not
identical. Nor is this source the separate positive dyadic theorem stream
proposed in the unary report. The retained calibration identifies
$\mathcal D_U\vdash1$ with a finite PA refutation, provably in PA.

\subsection{The actual blank startup and the protected tail}

The startup does not reach an ordinary clean source root. After
$48{,}140$ literal transitions, the final table reaches native state $84$
with head at $-2$, register header at $54$, source ordinal $136$, and
source address $001100110101010$. In register order
$(N,T,W,A,D,H,X,Y,L)$ its logical store is
\[
                         (1,0,13,0,0,1,0,0,1).
\]
The complete physical frame, including separators, spare registers,
address, and zero exterior, agrees with the retained startup frame.
The startup audit checks this literal assertion separately from the
unbounded runtime contracts.

The inherited suffix proof identifies this continuation with $141$
source primitive steps from ordinal $106$ and store
$(1,0,4,0,0,1,0,0,1)$. The incoming cache $1$ is an arithmetic selector
parameter, not an asserted theorem. The remaining axiom-builder suffix
creates nine explicit theorem heads above opaque numerical tail $4$:
seven logical instances, one dominated propositional instance, and truth
$0$. It returns to the scheduler with next candidate $1$, empty residual
pointer, theorem-valued top and cache, and cleared scratch registers.

The tail is unsafe if exposed: for Cantor pairing $P$,
$4=P(1,1)$ and $1=P(1,0)$. Its protection is therefore substantive.
The retained scheduler proof assigns conservative weights $-1$ to
selection, $0$ to modus ponens, and $+10$ to a completed ordinary axiom
command. Starting at candidate $1$, every command prefix incurs debt at
most two. The nine heads leave permanent margin seven. An explicit ghost
list length handles the alias $P(0,0)=0$; no numerical stack-length
decoder is assumed. Thus the tail never becomes exposed to top or cache,
and any halt supplies a numerical refutation.

Conversely, the retained strong compiler works above arbitrary underlying
stacks after resetting counter and cache. Successful candidate indices
are cofinal, and every unsuccessful candidate computation terminates.
The protected tail and finite initial cutoff consequently preserve
completeness as well as soundness~\cite{combined}.

\subsection{What the new reduction adds, and its scope}

The reassociations preserve every primitive effect and resolved ordinal
control-flow edge of the $3{,}909$-primitive source. The source-level
proof search and its calibration therefore remain the same. Connecting
the final literal table to that source additionally requires the
whole-tape primitive simulation contracts and the exact startup above.
The runtime verification described elsewhere in this report supplies
the final-table contracts; a source checksum or a finite startup alone
would not suffice.

Source induction, these runtime contracts, and the common blank frame
give the required blank-tape equivalence across the reassociations.
No arbitrary-tape equivalence is claimed for that address-changing step.
The separately proved predecessor-tag fold has its stronger all-tape
scope under the stated start-state and incoming-row hypotheses.

Composing the new arguments with the retained finite-refutation proofs
gives the ordinary mathematical conclusion
\[
 \PA\vdash
 \neg\Halt(M_{211},\mathrm{blank})
       \leftrightarrow\Con(\UU)
       \leftrightarrow\Con(\PA).
\]
The adjective ``inherited'' records proof provenance; it is not an added
object-theory assumption that PA is consistent. Unconditional nonhalting
is not asserted. The composition uses finite syntax, bounded local paths,
and natural-number induction, so it is PA-formalizable. Executing the
Python checkers is not itself checking an exported PA derivation: this
package contains no new end-to-end Lean, Rocq, or MM0 kernel certificate
for this literal machine.

\section{Rejected shortcuts and research implications}
\label{sec:research}

Several plausible reductions fail for concrete reasons. They are useful
because they explain why a smaller-looking source or a few successful
executions would not have been sufficient evidence for U211.

\paragraph{The original U212 patches are duplicates.}
All seven intact final tables are identical. Their local proofs can be
combined as redundant evidence, but their one-state reductions cannot be
counted repeatedly. The earlier known tagged rotations additionally
repurpose helper 211, invalidating the U212 bypass. Extending the closed
region to 936 primitives is what makes a different, compatible allocation
available here.

\paragraph{A shorter origin detector is not sound.}
The inherited source contains legal addresses with six consecutive zeros.
The supplied negative example shows a proposed six-zero detector taking
the wrong action on a genuine decrement frame. The final U211 source
retains maximum zero run six and the seven-zero detector. Early speculative
reassociations in this investigation produced runs of seven or eight zeros
and were discarded before being treated as valid machines.

\paragraph{Forward-prefix correctness is not the whole runtime.}
Another exploratory 211-state allocation preserved source order and obtained
the required address bound, but changed native 213 and 92. The literal row
$69/0=\act{0}{L}{213}$ reenters that part of the decoder during a primitive.
The full symbolic checker rejected the candidate at the zero case of a
decrement of register 8 at source ordinal 158. This counterexample motivated
the explicit backward-entry guard. The successful recipe does not change
213 or 92, and passes the full primitive replay. A source-tree theorem
cannot silently be upgraded to a theorem about the binary runtime.

\paragraph{Whole-state aliases did not suffice.}
A bounded screen considered all 44,310 ordered direct aliases obtained
by deleting a noninitial, nondispatch working state of U212 and retaining
another working state. It screened aliases against a suite containing two
small valuations of each of the 5,084 primitive cases, with a finite execution
cutoff; none passed all samples.
This is only a negative search result for that test and transformation
class. It is not a lower bound. The accepted construction changes the
source allocation first and uses a predecessor invariant instead of
requiring the two states to have identical behavior without a guard.

\paragraph{Theory size and machine size are different quantities.}
One schema need not compile into fewer states than a stronger theory with
an especially convenient encoding. The constant factors include the
formula syntax, arithmetic pairing, proof enumeration, decoder sharing,
and the literal initializer. Here the logical calculus and even the
197-node source DAG are unchanged; a single bit already written by a
predecessor saves one physical control state. The apparent arithmetic/set
theory reversal therefore does not imply a reversal of proof-theoretic
strength. This construction is an explicit upper bound at the stated
proof status; it makes no claim that 211 is optimal.

\section{Reproduction, provenance, and remaining formalization}
\label{sec:reproduction}

\subsection{Offline reproduction}

Extract the archive, enter its top-level directory, and run
\begin{lstlisting}
python3 reproduce.py
\end{lstlisting}
Python 3.10 or later and its standard library suffice. The default command
checks the frozen SHA-256 manifest, rebuilds both literal tables and the
source, independently audits all nine rotations, verifies the final fold
twice (including Python \code{-O}), checks the exact blank prefix, and
replays the complete symbolic primitive certificate. The portable runtime
replay took about 135 seconds in the preparation environment; machine speed
will change this. No network, discovery search, C++ compiler, or original
checkout is needed for the new U212-to-U211 verification.

\code{python3 reproduce.py --quick} skips the unbounded primitive replay and
reports that omission. It is a useful construction/local-proof check, not
an alternative proof of the full runtime. The symbolic engine contains
assertions and deliberately rejects \code{-O} or \code{PYTHONOPTIMIZE};
the full command refuses those settings. Its 4,000-inference per-case cap
is a proof-search limit, not a bound on the values certified by a success.

With a normal TeX installation, \code{python3 reproduce.py --quick --pdf}
also regenerates the appendix and compiles \path{report.tex} twice.
The paper uses standard \code{pdflatex} packages. The generated PDF's bytes
need not be reproducible across TeX versions or build timestamps; the
transition tables and source bytes are exact fixed outputs.

The integrity manifest covers frozen construction inputs, code, literal
machine files, paper sources and reference archives. The generated PDF and
replay outputs under \path{evidence/} are not sealed by that manifest,
because rebuilding rewrites them and timestamps or wall times vary. Their machine
identities are checked by the programs and recorded within the reports.
The archive itself provides the original evidence snapshots.

\subsection{What is supplied}

\begin{center}\small
\begin{tabularx}{\linewidth}{@{}p{38mm}X@{}}
\toprule
Path & Purpose\\
\midrule
\path{report.tex}, \path{paper/} & Complete mathematical report sources,
including the complete dense transition appendix.\\
\path{report.pdf} & Compiled, rendered and visually inspected report.\\
\path{machine/} & Native and dense U211 tables, rotated U212 intermediate,
new source sites, state isomorphism, and machine census.\\
\path{inputs/} & Frozen U212 parent/source, nine rotations and occurrence
paths, inherited label map, source manifest and damaged-archive recovery ledger.\\
\path{code/} & Deterministic builder, independent source/startup auditors,
local fold verifier, inherited symbolic engine and portable replay driver.\\
\path{evidence/} & Source, startup, fold and runtime results, all primitive-case
records, comparison manifest, and theory/runtime audit notes.\\
\path{reference/} & Complete compact U212 ancestry/calibration packet,
complete unary-theory packet, and preserved notices/licenses.\\
\bottomrule
\end{tabularx}
\end{center}

The new source and local machine arguments can be reviewed without accepting
the logical calibration. Conversely, the logical calibration cannot replace
the new runtime theorem. The dependency order is: inherited numerical
finite-refutation calibration; unchanged ordinal source behavior; final
literal primitive simulation; exact blank entry. Section~\ref{sec:calibration}
identifies the delicate protected-tail and compiler facts retained from the
baseline. The included reference packet makes those arguments inspectable
offline, but the new replay command does not claim to recheck a kernel
formalization of them.

\subsection{Exact output identities}

\begin{samepage}
The native U211 table SHA-256 is
\begin{center}\small
\path{7edae566b4c2e5c34d88030e0ae7caa49edda32c174b06b1a68bf9eb78567c8d}.
\end{center}
\end{samepage}
\begin{samepage}
The dense U211 table SHA-256 is
\begin{center}\small
\path{a529d2ea12b0f4c4185d2fe8ce9dd0f1633c04f39f6bf272bb37bf97bb31a349}.
\end{center}
\end{samepage}
\begin{samepage}
The reassociated source SHA-256 is
\begin{center}\small
\path{a1c81a3abff846a329a109927ce92d0be1553b7af2671855cdfe5bc4d6d953d1}.
\end{center}
\end{samepage}
\begin{samepage}
The rotated 212-state intermediate SHA-256 is
\begin{center}\small
\path{e3fe20c1d0181997354d0cf97b8ccceb85a1ec8055ad32ea5d02cacdf1a1b9d5}.
\end{center}
\end{samepage}
Hashes identify reviewed bytes; they do not prove semantic properties.
Both the source and runtime checks consume these literal outputs.

\subsection{Remaining proof-assistant work}

An end-to-end kernel result would need to formalize or import the exact
numerical-calculus and U/PA proof translations, the retained compiler and
startup-suffix theorem, the source relation, and the whole-tape symbolic
rules, then instantiate them on the released table. The local predecessor-tag
lemma is particularly small and could be formalized independently. No
missing kernel proof is represented here as an already completed artifact.
The present deliverable meets the requested below-212 candidate goal with
an exact 211-state table, a mathematical reduction, and reproducible
inductive verification at the stated boundary.


\begin{thebibliography}{9}
\bibitem{snapshot} ProveIt incoming research directory, fixed commit
\path{9d50bf3df8c9e14941b549945c4f8400172d81f4}.
\href{https://github.com/VladimirReshetnikov/ProveIt/tree/9d50bf3df8c9e14941b549945c4f8400172d81f4/docs/incoming}{Pinned source directory}.
Exact file URLs, sizes, Git blob identities and archive SHA-256 values are in
\path{inputs/source_manifest.json} and \path{evidence/input_comparison.json}.

\bibitem{unary} \emph{Unary power covering at PA strength}, 27 October 2026.
\path{pa_unary_power_cover.zip}, at the fixed snapshot~\cite{snapshot}.
Included intact in \path{reference/pa_unary_power_cover.zip}.
Relevant sections: unary core, unary arithmetic, upper transport, positive interface.

\bibitem{proposal1} \path{U212_ConPA_Stacked_2026-10-28 (1).zip},
at the fixed snapshot~\cite{snapshot}. Source of the inherited, unmodified
\path{symbolic_runtime.py}, \path{independent_runtime.py}, and legacy label map.
The other numbered proposals are identified in Table~\ref{tab:comparison}.

\bibitem{combined} \path{u212_combined_2026-10-29.zip},
at the fixed snapshot~\cite{snapshot}. Included intact in \path{reference/}.
Retains the direct U213 parent, U214/U218/U222 calibration and proofs,
literal-source and compiler arguments, startup suffix, and comparison inputs.

\bibitem{k215pa} Elliot Glazer, \emph{k215-pa}.
\url{https://github.com/elliotglazer/k215-pa}.
The fixed K215/ZFC theorem is background, not a proof certificate for U211.

\bibitem{k230} Elliot Glazer, \emph{k230-rocq}, ``The theory K.''
\url{https://github.com/elliotglazer/k230-rocq#the-theory-k}.

\bibitem{m254} \emph{busy-beaver}, branch \path{zf-254}, \path{ZF/m254}.
\url{https://github.com/jeff87654/busy-beaver/tree/zf-254/ZF/m254}.
Historical machine and compiler context supplied with the research question.
\end{thebibliography}

\clearpage
\appendix
\section{Complete dense transition table}
\label{sec:table}
States are $0,\ldots,210$; $H$ denotes the separate halt label 211.
The initial state is 0, the head starts at coordinate 0, and every tape cell
is initially 0. A row applies its write and movement before entering its
next state. All moves are one cell left or right. There are no implicit
transitions or uncounted preparatory states. This appendix is generated
directly from \path{machine/u211.tm}; its three groups of columns list the
states of one machine.
\begingroup\footnotesize
\setlength{\tabcolsep}{6pt}\renewcommand{\arraystretch}{1.08}
\begin{longtable}{rll@{\hspace{5mm}}rll@{\hspace{5mm}}rll}
\caption{The complete dense U211 table. Each action is write, move, next state.}\label{tab:machine}\\
\toprule State & Read 0 & Read 1 & State & Read 0 & Read 1 & State & Read 0 & Read 1\\\midrule
\endfirsthead
\toprule State & Read 0 & Read 1 & State & Read 0 & Read 1 & State & Read 0 & Read 1\\\midrule
\endhead
\midrule\multicolumn{9}{r}{Continued on next page}\\\endfoot
\bottomrule\endlastfoot
0 & $1\,R\,92$ & $1\,R\,0$ & 71 & $0\,R\,55$ & $1\,L\,79$ & 142 & $0\,R\,12$ & $1\,R\,159$\\
1 & $0\,L\,1$ & $1\,L\,2$ & 72 & $0\,R\,80$ & $1\,R\,81$ & 143 & $0\,R\,88$ & $1\,L\,166$\\
2 & $0\,L\,3$ & $1\,L\,2$ & 73 & $1\,R\,0$ & $1\,R\,63$ & 144 & $0\,R\,88$ & $1\,R\,158$\\
3 & $0\,L\,5$ & $1\,L\,2$ & 74 & $0\,R\,74$ & $1\,R\,82$ & 145 & $0\,L\,152$ & $1\,L\,163$\\
4 & $1\,R\,6$ & $0\,L\,4$ & 75 & $0\,L\,145$ & $1\,L\,141$ & 146 & $0\,R\,84$ & $1\,R\,88$\\
5 & $0\,L\,9$ & $1\,L\,2$ & 76 & $0\,R\,85$ & $1\,R\,86$ & 147 & $0\,R\,159$ & $1\,R\,12$\\
6 & $0\,R\,6$ & $0\,R\,10$ & 77 & $0\,R\,20$ & $1\,R\,38$ & 148 & $0\,R\,83$ & $1\,R\,143$\\
7 & $0\,R\,84$ & $1\,R\,191$ & 78 & $0\,R\,89$ & $1\,L\,132$ & 149 & $0\,R\,144$ & $1\,R\,155$\\
8 & $0\,R\,14$ & $1\,R\,13$ & 79 & $0\,L\,108$ & $1\,R\,56$ & 150 & $0\,L\,66$ & $1\,L\,141$\\
9 & $0\,L\,15$ & $1\,L\,2$ & 80 & $0\,R\,49$ & $0\,R\,33$ & 151 & $0\,R\,203$ & $1\,R\,170$\\
10 & $0\,L\,16$ & $0\,R\,10$ & 81 & $0\,R\,91$ & $0\,L\,42$ & 152 & $0\,R\,159$ & $1\,R\,114$\\
11 & $0\,L\,26$ & $1\,R\,128$ & 82 & $0\,R\,179$ & $1\,R\,77$ & 153 & $0\,R\,97$ & $1\,R\,97$\\
12 & $0\,R\,19$ & $1\,R\,20$ & 83 & $0\,R\,96$ & $1\,R\,12$ & 154 & $0\,R\,49$ & $0\,R\,32$\\
13 & $0\,L\,89$ & $1\,L\,31$ & 84 & $0\,R\,88$ & $1\,R\,97$ & 155 & $0\,R\,189$ & $1\,R\,112$\\
14 & $0\,L\,11$ & $0\,R\,22$ & 85 & $0\,R\,209$ & $1\,R\,21$ & 156 & $0\,R\,196$ & $1\,R\,203$\\
15 & $0\,L\,23$ & $1\,L\,2$ & 86 & $0\,R\,116$ & $1\,R\,100$ & 157 & $0\,L\,188$ & $1\,R\,96$\\
16 & $1\,L\,1$ & $0\,L\,1$ & 87 & $0\,R\,148$ & $1\,R\,167$ & 158 & $0\,R\,168$ & $1\,R\,180$\\
17 & $0\,R\,24$ & $0\,R\,32$ & 88 & $0\,R\,102$ & $1\,R\,103$ & 159 & $0\,R\,19$ & $1\,R\,171$\\
18 & $0\,R\,98$ & $1\,R\,113$ & 89 & $0\,R\,80$ & $1\,R\,79$ & 160 & $0\,R\,164$ & $1\,R\,188$\\
19 & $0\,R\,28$ & $1\,R\,29$ & 90 & $0\,R\,101$ & $1\,R\,197$ & 161 & $0\,R\,159$ & $1\,R\,96$\\
20 & $0\,R\,30$ & $1\,R\,31$ & 91 & $0\,R\,105$ & $1\,R\,106$ & 162 & $0\,R\,125$ & $1\,R\,164$\\
21 & $0\,R\,24$ & $0\,R\,59$ & 92 & $0\,R\,65$ & $0\,R\,0$ & 163 & $0\,R\,97$ & $1\,R\,162$\\
22 & $0\,R\,33$ & $1\,R\,22$ & 93 & $0\,R\,27$ & $1\,R\,107$ & 164 & $0\,L\,7$ & $1\,R\,142$\\
23 & $0\,L\,34$ & $1\,L\,2$ & 94 & $0\,R\,184$ & $0\,L\,210$ & 165 & $0\,L\,66$ & $1\,R\,174$\\
24 & $0\,L\,35$ & $0\,R\,0$ & 95 & $1\,L\,109$ & $1\,L\,95$ & 166 & $0\,R\,146$ & $1\,R\,173$\\
25 & $0\,R\,36$ & $1\,R\,25$ & 96 & $0\,R\,19$ & $1\,R\,110$ & 167 & $0\,R\,175$ & $1\,R\,202$\\
26 & $0\,R\,37$ & $1\,R\,38$ & 97 & $0\,R\,60$ & $1\,R\,60$ & 168 & $0\,R\,161$ & $1\,R\,143$\\
27 & $0\,R\,39$ & $0\,L\,40$ & 98 & $0\,R\,68$ & $1\,R\,135$ & 169 & $0\,R\,126$ & $1\,R\,87$\\
28 & $0\,R\,41$ & $1\,R\,42$ & 99 & $0\,R\,93$ & $1\,R\,108$ & 170 & $0\,R\,161$ & $1\,R\,96$\\
29 & $0\,R\,43$ & $1\,R\,44$ & 100 & $0\,R\,105$ & $0\,L\,210$ & 171 & $0\,R\,182$ & $1\,R\,122$\\
30 & $0\,R\,45$ & $1\,R\,46$ & 101 & $0\,R\,117$ & $1\,R\,118$ & 172 & $0\,R\,145$ & $1\,R\,196$\\
31 & $0\,R\,21$ & $1\,R\,47$ & 102 & $0\,R\,71$ & $1\,R\,58$ & 173 & $0\,L\,194$ & $1\,R\,60$\\
32 & $0\,R\,48$ & $1\,R\,32$ & 103 & $0\,R\,57$ & $1\,R\,119$ & 174 & $0\,R\,145$ & $1\,R\,88$\\
33 & $0\,R\,32$ & $1\,R\,33$ & 104 & $0\,R\,21$ & $1\,R\,120$ & 175 & $0\,R\,152$ & $1\,R\,204$\\
34 & $0\,L\,74$ & $1\,L\,2$ & 105 & $0\,R\,21$ & $0\,L\,10$ & 176 & $0\,R\,140$ & $1\,R\,84$\\
35 & $1\,R\,35$ & $1\,L\,1$ & 106 & $0\,R\,61$ & $1\,R\,42$ & 177 & $0\,R\,176$ & $1\,R\,75$\\
36 & $0\,R\,50$ & $1\,R\,36$ & 107 & $0\,R\,121$ & $1\,R\,8$ & 178 & $0\,R\,146$ & $1\,R\,169$\\
37 & $0\,R\,17$ & $1\,R\,111$ & 108 & $0\,R\,101$ & $1\,R\,122$ & 179 & $0\,R\,138$ & $1\,R\,116$\\
38 & $0\,R\,13$ & $1\,R\,52$ & 109 & $0\,L\,210$ & $0\,L\,95$ & 180 & $0\,R\,97$ & $1\,R\,51$\\
39 & $0\,L\,11$ & $0\,R\,36$ & 110 & $0\,R\,124$ & $1\,R\,104$ & 181 & $0\,R\,168$ & $1\,R\,51$\\
40 & $0\,L\,4$ & $1\,L\,53$ & 111 & $0\,R\,115$ & $1\,L\,72$ & 182 & $0\,R\,17$ & $1\,L\,131$\\
41 & $0\,R\,24$ & $0\,R\,25$ & 112 & $0\,R\,165$ & $1\,R\,182$ & 183 & $0\,R\,143$ & $1\,R\,149$\\
42 & $0\,R\,13$ & $0\,L\,16$ & 113 & $0\,R\,160$ & $1\,R\,185$ & 184 & $0\,R\,76$ & $1\,R\,116$\\
43 & $0\,R\,55$ & $1\,R\,56$ & 114 & $0\,R\,148$ & $1\,R\,146$ & 185 & $0\,R\,106$ & $1\,R\,124$\\
44 & $0\,R\,57$ & $1\,R\,58$ & 115 & $0\,R\,11$ & $1\,R\,129$ & 186 & $1\,L\,40$ & $1\,L\,109$\\
45 & $0\,R\,24$ & $0\,R\,54$ & 116 & $0\,L\,131$ & $0\,L\,1$ & 187 & $0\,R\,172$ & $1\,R\,159$\\
46 & $0\,R\,60$ & $1\,R\,61$ & 117 & $0\,R\,39$ & $1\,R\,72$ & 188 & $0\,R\,183$ & $1\,R\,184$\\
47 & $0\,R\,62$ & $0\,L\,16$ & 118 & $0\,R\,61$ & $1\,R\,132$ & 189 & $0\,R\,193$ & $1\,R\,207$\\
48 & $0\,R\,25$ & $1\,R\,48$ & 119 & $0\,R\,61$ & $1\,R\,133$ & 190 & $0\,R\,140$ & $1\,R\,199$\\
49 & $0\,L\,35$ & $0\,R\,63$ & 120 & $0\,R\,70$ & $0\,L\,16$ & 191 & $0\,R\,157$ & $1\,R\,192$\\
50 & $0\,R\,16$ & $1\,R\,50$ & 121 & $0\,R\,14$ & $1\,R\,132$ & 192 & $0\,R\,176$ & $1\,R\,187$\\
51 & $0\,R\,147$ & $1\,R\,143$ & 122 & $0\,R\,21$ & $1\,R\,134$ & 193 & $0\,R\,64$ & $1\,R\,7$\\
52 & $0\,R\,67$ & $1\,R\,68$ & 123 & $0\,R\,151$ & $1\,R\,190$ & 194 & $0\,R\,150$ & $1\,R\,127$\\
53 & $0\,L\,40$ & $1\,L\,53$ & 124 & $0\,R\,137$ & $1\,L\,20$ & 195 & $0\,R\,18$ & $1\,R\,94$\\
54 & $0\,R\,22$ & $1\,R\,54$ & 125 & $0\,R\,130$ & $1\,R\,181$ & 196 & $0\,R\,203$ & $1\,R\,174$\\
55 & $0\,L\,11$ & $0\,R\,50$ & 126 & $0\,R\,206$ & $1\,R\,172$ & 197 & $0\,R\,104$ & $1\,R\,96$\\
56 & $0\,R\,60$ & $0\,L\,16$ & 127 & $0\,R\,75$ & $1\,R\,177$ & 198 & $0\,R\,90$ & $1\,R\,88$\\
57 & $0\,R\,69$ & $1\,R\,70$ & 128 & $0\,R\,75$ & $1\,R\,208$ & 199 & $0\,R\,139$ & $1\,R\,201$\\
58 & $0\,R\,61$ & $1\,R\,71$ & 129 & $0\,R\,205$ & $1\,R\,157$ & 200 & $0\,R\,96$ & $1\,R\,123$\\
59 & $0\,R\,54$ & $1\,R\,59$ & 130 & $0\,R\,12$ & $1\,R\,141$ & 201 & $0\,R\,204$ & $1\,R\,88$\\
60 & $0\,R\,49$ & $0\,R\,25$ & 131 & $0\,R\,30$ & $1\,R\,69$ & 202 & $0\,R\,200$ & $1\,R\,156$\\
61 & $0\,L\,28$ & $1\,L\,72$ & 132 & $0\,R\,39$ & $1\,L\,124$ & 203 & $0\,R\,88$ & $1\,R\,84$\\
62 & $0\,L\,117$ & $1\,L\,30$ & 133 & $0\,R\,153$ & $0\,L\,16$ & 204 & $0\,L\,123$ & $1\,R\,168$\\
63 & $0\,R\,73$ & $1\,R\,63$ & 134 & $0\,R\,154$ & $0\,L\,16$ & 205 & $0\,R\,198$ & $1\,R\,28$\\
64 & $0\,R\,97$ & $1\,R\,150$ & 135 & $0\,R\,88$ & $1\,R\,136$ & 206 & $0\,R\,163$ & $1\,R\,194$\\
65 & $1\,R\,186$ & $1\,L\,53$ & 136 & $0\,R\,198$ & $1\,R\,140$ & 207 & $0\,R\,66$ & $1\,R\,178$\\
66 & $0\,R\,144$ & $1\,R\,145$ & 137 & $0\,R\,24$ & $0\,R\,48$ & 208 & $0\,R\,161$ & $1\,R\,195$\\
67 & $0\,R\,76$ & $1\,R\,H$ & 138 & $0\,R\,99$ & $1\,R\,37$ & 209 & $0\,R\,45$ & $1\,R\,100$\\
68 & $0\,L\,198$ & $1\,R\,31$ & 139 & $0\,R\,173$ & $1\,R\,148$ & 210 & $0\,L\,4$ & $0\,L\,210$\\
69 & $0\,R\,41$ & $1\,R\,62$ & 140 & $0\,R\,96$ & $1\,R\,159$ &  &  & \\
70 & $0\,L\,89$ & $1\,R\,78$ & 141 & $0\,R\,176$ & $1\,R\,166$ &  &  & \\
\end{longtable}
\endgroup

\end{document}
